% !TITLE 浣溪沙·一曲新词酒一杯
% !AUTHOR 晏殊
% !DYNASTY 两宋
% !GENRE 词
% !ORDER 1
% !EDITION 《全宋词》，唐圭璋编，中华书局
% !EDITION_STATUS pending

\begin{poem}
一曲新词酒一杯，去年天气旧亭台\textsuperscript{1}。夕阳西下几时回\textsuperscript{2}？

无可奈何\textsuperscript{3}花落去，似曾相识\textsuperscript{4}燕归来。小园香径\textsuperscript{5}独徘徊\textsuperscript{6}。
\end{poem}

\begin{annotations}
\notetext{1}{旧亭台：去年的亭台楼阁，和去年一模一样。“去年天气旧亭台”——天气还是去年的天气，亭台还是从前的亭台，但人已不同。}
\notetext{2}{几时回：什么时候回来？夕阳落下去了，但落下的夕阳还能再升起，逝去的时光呢？}
\notetext{3}{无可奈何：没有办法。花落是自然规律，谁也阻止不了——这句话后来成了成语。}
\notetext{4}{似曾相识：好像曾经见过。燕子回来了，像是去年的那只——但“似曾”二字说明并不确定，物是人非。}
\notetext{5}{香径：落花飘香的小路。园中小路因落花而带香气。}
\notetext{6}{徘徊：走来走去，流连不去。独自一人在落花飘香的小路上踱步，写尽了孤独与无奈。}
\end{annotations}

\begin{appreciation}
《浣溪沙》是晏殊词中最负盛名的一首，也是北宋士大夫闲愁词的代表之一。晏殊十四岁以神童入试，官至宰相，一生顺遂无波，可偏偏在最不该发愁的处境里，写出了对时光最无可奈何的一声叹息。

“一曲新词酒一杯，去年天气旧亭台。”新谱的曲子，新斟的酒，天气还是去年的天气，亭台还是从前的亭台。一切照旧，可坐在其中的人知道又一年已经过去了。“夕阳西下几时回？”太阳落了还会升起，今天的落日却不再是昨天那一轮；这一问问的是落日，问的其实是去而不返的光阴。

下阕两句是全词的眼：“无可奈何花落去，似曾相识燕归来。”花要落，任凭你多少权势手段也终究拦不住；燕子归来，像极了去年那一只，却又不敢认实。一去一来，一无可挽回，一仿佛相识，时间的两副面孔在这一联里对得严丝合缝。这一联又在晏殊的七律《示张寺丞王校勘》中出现，诗词并用，可见他自己也得意。宋人笔记还有一个更美的传闻，说晏殊只得了上句，苦思无以对，是江都尉王琪替他补出下句——此事未必可信，却足见后人对这一联的珍爱。

“小园香径独徘徊。”落花铺地，香气满径，一个人在那里来回地走。全词没有一句呼天抢地，末了只落在一个“独”字上——府中宾客盈门，可黄昏的小径上，谁都只有自己。

《生年不满百》劝人“为乐当及时”，是把忧患推开；晏殊不推开，他承认它。花落就花落，燕来就燕来，人所能做的，不过是站在香径上多看一会儿。富贵买不通时间，权势也留不住春天，这一点上，宰相与平民并没有两样。
\end{appreciation}
